
% ==============================================================
% סמינר בהתמרת לפלס | מצגת Beamer בעברית
% Compiler: XeLaTeX | Beamer 16:9
% Manim animations are separate MP4s.
% ==============================================================
\documentclass[aspectratio=169,11pt]{beamer}

% Mathematical and graphics packages - load BEFORE polyglossia
\usepackage{amsmath,amssymb,mathtools}
\usepackage{graphicx}
\usepackage{tikz}
\usepackage{hyperref}

% Hebrew and English support
\usepackage{fontspec}
\usepackage{polyglossia}
\setdefaultlanguage{hebrew}
\setotherlanguage{english}

\newfontfamily\hebrewfont[Script=Hebrew]{FreeSans}
\newfontfamily\hebrewfontsf[Script=Hebrew]{FreeSans}
\newfontfamily\hebrewfonttt[Script=Hebrew]{FreeMono}
\setsansfont{FreeSans}

\usefonttheme{professionalfonts}
\setbeamersize{text margin left=11mm,text margin right=11mm}


% Colors
\definecolor{night}{HTML}{0B1428}
\definecolor{panel}{HTML}{14223C}
\definecolor{lighttext}{HTML}{F1F6FC}
\definecolor{mutedtext}{HTML}{B6C5D8}
\definecolor{cyan}{HTML}{41C9E8}
\definecolor{gold}{HTML}{FFCD66}
\definecolor{rose}{HTML}{FB7185}
\definecolor{mint}{HTML}{7BDCB5}
\setbeamercolor{normal text}{fg=lighttext,bg=night}
\setbeamercolor{structure}{fg=cyan}
\setbeamercolor{alerted text}{fg=gold}
\setbeamercolor{frametitle}{fg=lighttext,bg=night}
\setbeamercolor{block title}{fg=night,bg=cyan}
\setbeamercolor{block body}{fg=lighttext,bg=panel}
\setbeamertemplate{navigation symbols}{}
\setbeamertemplate{itemize item}{\textcolor{cyan}{\small$\blacktriangleright$}}
\setbeamertemplate{itemize subitem}{\textcolor{gold}{\scriptsize$\bullet$}}
\setbeamertemplate{frametitle}{%
\vspace{2mm}\begin{flushright}{\Large\bfseries\insertframetitle}\end{flushright}%
\vspace{-2.7mm}\textcolor{cyan}{\rule{\linewidth}{0.7pt}}\vspace{0.8mm}}
\setbeamertemplate{footline}{%
  \hspace{7mm}{\tiny\color{mutedtext}סמינר: התמרת לפלס}\hfill
  {\tiny\color{mutedtext}\insertframenumber}\hspace{7mm}\vspace{3mm}}
\setlength{\parskip}{5pt}
\newcommand{\Lap}{\mathcal{L}}
\newcommand{\Real}{\operatorname{Re}}
\newcommand{\dd}{\,\mathrm{d}}
\newcommand{\highlight}[1]{{\color{gold}#1}}
\newcommand{\bigformula}[1]{\vspace{2mm}\begin{center}{\LARGE\color{cyan}\(\displaystyle #1\)}\end{center}\vspace{2mm}}
\newcommand{\visualframe}[3]{%
\begin{frame}{#1}
\begin{center}
\includegraphics[width=.96\linewidth,height=.74\textheight,keepaspectratio]{assets/#2}
\end{center}
\vspace{-1mm}\begin{center}{\small\color{mutedtext}#3}\end{center}
\end{frame}}
\newcommand{\twoline}[2]{\begin{block}{#1}\vspace{1mm}#2\vspace{1mm}\end{block}}
\newenvironment{spaceditems}{\begin{itemize}\setlength{\itemsep}{8pt}}{\end{itemize}}

\title{התמרת לפלס: מן האקספוננט לפתרון בעיות התחלה}
\subtitle{הגדרה, משפטים, הוכחות ויישומים}
\author{רועי אוחנה}
\date{2026}

\begin{document}

% 01
\begin{frame}[plain]
\begin{center}
\vspace{16mm}
{\color{cyan}\rule{.26\linewidth}{1.2pt}}\\[7mm]
{\Huge\bfseries התמרת לפלס}\\[4mm]
{\Large מן האקספוננט לפתרון בעיות התחלה}\\[10mm]
{\large\color{mutedtext}סמינר במתמטיקה \quad | \quad רועי אוחנה}\\[13mm]
{\color{gold}\LARGE $\displaystyle F(s)=\int_0^\infty e^{-st}f(t)\dd t$}
\end{center}
\end{frame}

% 02
\begin{frame}{מפת הדרך להרצאה}
\begin{spaceditems}
\item \highlight{מוטיבציה:} מדוע האקספוננט הוא רכיב טבעי של משוואות דיפרנציאליות?
\item \highlight{יסודות:} הגדרת ההתמרה, התכנסות ותכונות מרכזיות
\item \highlight{הכלי המרכזי:} המרת משוואות דיפרנציאליות למשוואות אלגבריות
\item \highlight{קלטים מיוחדים:} השהיה בזמן ודלתא של דיראק
\item \highlight{מערכות לינאריות:} קונבולוציה, תגובת אימפולס ומצב מתמיד
\end{spaceditems}
\vspace{6mm}
\begin{center}{\color{cyan}פונקציה בזמן $\longleftarrow$ ייצוג בתחום $s$ $\longleftarrow$ פתרון וחזרה לזמן}\end{center}
\end{frame}

% 03
\begin{frame}{המוטיבציה: האקספוננט תחת גזירה}
\begin{spaceditems}
\item הפונקציה האקספוננציאלית שומרת על צורתה גם לאחר גזירה
\item לכן היא מועמדת טבעית לפתרון משוואות לינאריות עם מקדמים קבועים
\end{spaceditems}
\bigformula{\frac{\mathrm d^n}{\mathrm dt^n} e^{st}=s^n e^{st}}
למשל, אם $ay''+by'+cy=0$ וננסה $y(t)=e^{st}$, נקבל
\bigformula{as^2+bs+c=0}
\begin{center}{\small\color{mutedtext}המשוואה הדיפרנציאלית הוחלפה במשוואה אלגברית.}\end{center}
\end{frame}

% 04
\visualframe{המחשה: נגזרת האקספוננט}{01_exponential.png}{המשיק משתנה בהתאם לגובה הגרף; ההמחשה מלווה גם בסרטון אנימציה.}

% 05
\begin{frame}{האקספוננט המרוכב ונוסחת אוילר}
\bigformula{s=\sigma+i\omega}
\bigformula{e^{st}=e^{\sigma t}(\cos\omega t+i\sin\omega t)}
\begin{spaceditems}
\item החלק הממשי $\sigma$ קובע גדילה או דעיכה אקספוננציאלית
\item החלק המדומה $\omega$ קובע תדירות תנודה
\item צירוף של מצבים כאלה מסביר מגוון תופעות דינמיות
\end{spaceditems}
\end{frame}

% 06
% ANIM: euler_decay
\visualframe{המחשה: סיבוב ודעיכה}{02_euler.png}{נוסחת אוילר מתארת סיבוב; הכפלה באקספוננט ממשי מתארת דעיכה או גדילה.}

% 07
\begin{frame}{מצבים אקספוננציאליים במישור המרוכב}
אפשר לפרש פתרונות של מערכות לינאריות כצירוף של מצבים:
\bigformula{x(t)=\sum_{k=1}^{n}c_k e^{s_k t}}
\begin{spaceditems}
\item לכל מספר $s_k$ יש חלק ממשי וחלק מדומה
\item מיקום $s_k$ במישור מתאר את קצב הדעיכה ואת תדירות התנודה
\item זהו מניע רעיוני להתמרת לפלס; ההתמרה עצמה אינה סתם פירוק לסכום אקספוננטים
\end{spaceditems}
\end{frame}

% 08
% ANIM: complex_modes
\visualframe{המחשה: מיקום במישור $s$}{03_modes.png}{מיקום הנקודות במישור מרוכב משנה את ההתנהגות בזמן.}

% 09
\begin{frame}{הגדרת התמרת לפלס}
\begin{block}{הגדרה}
בהינתן פונקציה $f:[0,\infty)\to\mathbb C$, נגדיר את התמרת לפלס החד־צדדית שלה בכל נקודה שבה האינטגרל מתכנס:
\end{block}
\bigformula{F(s)=\Lap\{f\}(s)=\int_0^\infty e^{-st}f(t)\dd t}
\begin{spaceditems}
\item $t$ הוא המשתנה המקורי; $s$ הוא משתנה מרוכב חדש
\item הגורם $e^{-st}$ משקלל את התרומה של $f$ לאורך הזמן
\item ההתמרה מוגדרת רק בתחום ההתכנסות של האינטגרל
\end{spaceditems}
\end{frame}

% 10
\visualframe{המחשה: שקלול הפונקציה באקספוננט}{04_weighted.png}{האינטגרל אוסף את התרומות המשוקללות של האות על פני הזמן.}

% 11
\begin{frame}{מתי התמרת לפלס קיימת?}
\begin{block}{סדר אקספוננציאלי}
אם קיימים $M>0$, $\alpha\in\mathbb R$ ו־$t_0\ge0$ כך שלכל $t\ge t_0$ מתקיים
\[|f(t)|\le Me^{\alpha t},\]
נאמר שלפונקציה יש סדר אקספוננציאלי $\alpha$.
\end{block}
\begin{block}{משפט קיום — תנאי מספיק}
אם $f$ רציפה למקוטעין על כל קטע סופי והיא מסדר אקספוננציאלי $\alpha$, אז
\[\highlight{\Lap\{f\}(s)\ \text{קיימת כאשר}\ \Real(s)>\alpha}.\]
\end{block}
\vspace{2mm}
{\small הרעיון בהוכחה: השוואת $|e^{-st}f(t)|$ ל־$Me^{-(\Real(s)-\alpha)t}$.}
\end{frame}

% 12
\visualframe{המחשה: תנאי הגדילה}{05_growth.png}{פונקציה כמו $e^{t^2}$ גדלה מהר יותר מכל $e^{\alpha t}$ עבור $\alpha$ קבוע.}

% 13
\begin{frame}{התמרות בסיסיות ולינאריות}
\begin{block}{זוגות התמרה שימושיים}
\begin{align*}
\Lap\{1\}&=\frac1s & \Lap\{\sin(\omega t)\}&=\frac{\omega}{s^2+\omega^2}\\[2mm]
\Lap\{e^{at}\}&=\frac1{s-a} & \Lap\{\cos(\omega t)\}&=\frac{s}{s^2+\omega^2}\\[2mm]
\Lap\{t^n\}&=\frac{n!}{s^{n+1}} &&
\end{align*}
\end{block}
\vspace{5mm}
\begin{center}{\large\color{gold}$\displaystyle\Lap\{af+bg\}=a\Lap\{f\}+b\Lap\{g\}$}\end{center}
\begin{center}{\small\color{mutedtext}כל זה בתחום שבו ההתמרות מתכנסות.}\end{center}
\end{frame}

% 14
\begin{frame}{משפט ההזזה הראשון}
\bigformula{\Lap\{e^{at}f(t)\}(s)=F(s-a)}
\begin{block}{רעיון ההוכחה}
\begin{align*}
\Lap\{e^{at}f(t)\}(s)
&=\int_0^\infty e^{-st}e^{at}f(t)\dd t\\
&=\int_0^\infty e^{-(s-a)t}f(t)\dd t\\
&=F(s-a).
\end{align*}
\end{block}
\begin{center}{\small\color{mutedtext}כפל באקספוננט בזמן פירושו הזזה ב־$s$.}\end{center}
\end{frame}

% 15
\visualframe{המחשה: הזזה במישור ההתמרה}{06_shift_s.png}{הכפלה ב־$e^{at}$ מזיזה את מבנה ההתמרה ב־$a$ לאורך החלק הממשי.}

% 16
\begin{frame}{למה נגזרות הופכות לביטויים אלגבריים?}
\bigformula{\Lap\{f'(t)\}=sF(s)-f(0)}
\begin{block}{הוכחה באינטגרציה בחלקים}
\begin{align*}
\int_0^\infty e^{-st}f'(t)\dd t
&=\left[e^{-st}f(t)\right]_0^\infty
+s\int_0^\infty e^{-st}f(t)\dd t\\
&=sF(s)-f(0).
\end{align*}
\end{block}
עבור נגזרת שנייה מתקבל
\[\highlight{\Lap\{f''(t)\}=s^2F(s)-sf(0)-f'(0)}.\]
\end{frame}

% 17
\begin{frame}{התמרת לפלס הפוכה ושברים חלקיים}
המטרה כעת היא לחזור מפונקציה של $s$ אל הפונקציה המקורית בזמן.
\[\frac{1}{(s-2)(s-3)}=-\frac{1}{s-2}+\frac{1}{s-3}\]
\bigformula{\Lap^{-1}\!\left\{\frac{1}{(s-2)(s-3)}\right\}
=-e^{2t}+e^{3t}}
\begin{spaceditems}
\item מפרקים פונקציה רציונלית לשברים חלקיים
\item משתמשים בזוגות ההתמרה הידועים, כולל משפטי ההזזה
\item תנאי הקיום והיחידות רלוונטיים כאשר מזהים את ההתמרה ההפוכה
\end{spaceditems}
\end{frame}

% 18
\begin{frame}{יישום ראשון: פתרון בעיית התחלה}
נפתור את המשוואה
\[y''+3y'+2y=0,\qquad y(0)=1,\quad y'(0)=0.\]
נסמן $Y(s)=\Lap\{y\}$. על פי נוסחאות הגזירה:
\begin{align*}
(s^2Y-s)+3(sY-1)+2Y&=0\\[1mm]
Y(s)&=\frac{s+3}{(s+1)(s+2)}\\
&=\frac{2}{s+1}-\frac{1}{s+2}.
\end{align*}
\bigformula{y(t)=2e^{-t}-e^{-2t}}
\end{frame}

% 19
\visualframe{המחשה: הפתרון בזמן}{07_ode.png}{הפתרון מקיים את תנאי ההתחלה ומתכנס ל־0 כאשר $t\to\infty$.}

% 20
\begin{frame}{קלט המופעל באיחור: משפט ההזזה השני}
נסמן את פונקציית המדרגה בזמן $a>0$ ב־$u_a(t)$:
\[u_a(t)=\begin{cases}0,&t<a,\\1,&t\ge a.\end{cases}\]
\bigformula{\Lap\{u_a(t)f(t-a)\}=e^{-as}F(s)}
\begin{block}{הוכחה קצרה}
מבצעים באינטגרל את שינוי המשתנה $\tau=t-a$:
\[\int_a^\infty e^{-st}f(t-a)\dd t
=e^{-as}\int_0^\infty e^{-s\tau}f(\tau)\dd\tau.\]
\end{block}
\end{frame}

% 21
% ANIM: time_shift
\visualframe{המחשה: אות מושהה בזמן}{08_time_shift.png}{האות נשאר כבוי עד זמן $a$, ואז מתחיל להתפתח כאילו הזמן אופס.}

% 22
\begin{frame}{דלתא של דיראק: מודל לאימפולס}
\begin{spaceditems}
\item כוח גדול שפועל זמן קצר יכול לתת שינוי תנע סופי
\item דלתא מתארת באופן אידיאלי אירוע מרוכז בזמן יחיד
\end{spaceditems}
\[\int_{-\infty}^\infty f(t)\delta(t-a)\dd t=f(a)\]
\bigformula{\Lap\{\delta(t-a)\}=e^{-as}\qquad(a>0)}
\begin{block}{דיוק מתמטי}
דלתא אינה פונקציה רגילה: אפשר להגדירה כהתפלגות, או לפרש אינטגרלים מתאימים באמצעות רימן–שטילטייס.
\end{block}
\end{frame}

% 23
% ANIM: dirac_pulse
\visualframe{המחשה: פולס צר כמקורב לדלתא}{09_dirac.png}{הרוחב קטן, הגובה גדל, אך השטח נשאר 1.}

% 24
\begin{frame}{דוגמה פיזיקלית: אימפולס ותנועה בגז צמיג}
קליע שמסתו $m$ מקבל אימפולס של תנע $mv_0$ ואז נע בתווך צמיג:
\[mx''(t)+kx'(t)=mv_0\delta(t).\]
נשתמש בקונבנציה החד־צדדית הסיבתית עם מצב מנוחה לפני האימפולס:
\[X(s)=\frac{v_0}{s(s+k/m)}.\]
\bigformula{x(t)=\frac{mv_0}{k}\left(1-e^{-kt/m}\right)}
\[x'(t)=v_0e^{-kt/m}\qquad (t>0).\]
\begin{center}{\small\color{mutedtext}האימפולס יוצר קפיצה מיידית במהירות; לאחר מכן ההתנגדות מאטה את הקליע.}\end{center}
\end{frame}

% 25
\visualframe{המחשה: מהירות ומיקום אחרי אימפולס}{10_impulse_motion.png}{המהירות דועכת, בעוד המיקום מתקרב לערך סופי.}

% 26
\begin{frame}{קונבולוציה ומשפט הקונבולוציה}
\begin{block}{הגדרה}
\[(f*g)(t)=\int_0^t f(\tau)g(t-\tau)\dd\tau.\]
\end{block}
\bigformula{\Lap\{f*g\}(s)=F(s)G(s)}
\begin{block}{רעיון ההוכחה}
\[F(s)G(s)=\int_0^\infty\!\int_0^\infty
 e^{-s(u+v)}f(u)g(v)\dd u\dd v.\]
בהנחות המאפשרות שינוי סדר אינטגרציה, מציבים $t=u+v$ ומתקבלת התמרת הקונבולוציה.
\end{block}
\end{frame}

% 27
% ANIM: convolution
\visualframe{המחשה: קונבולוציה כשטח חפיפה}{11_convolution.png}{כאשר אחת הפונקציות מחליקה, שטח החפיפה קובע את ערך הקונבולוציה.}

% 28
\begin{frame}{מתגובה לאימפולס לתגובה לקלט כללי}
במערכת לינארית ואינווריאנטית בזמן, בעלת תגובת אימפולס $h(t)$:
\bigformula{y(t)=(x*h)(t)}
ובתחום לפלס מתקיים
\bigformula{Y(s)=X(s)H(s)}
\begin{spaceditems}
\item תגובת אימפולס אחת מאפשרת לחשב תגובות לקלטים אחרים
\item זהו הקשר בין משפט הקונבולוציה לתורת המערכות
\item למשל, אם $h(t)=e^{-t}u_0(t)$ ו־$x(t)=u_1(t)$ אז
\end{spaceditems}
\[y(t)=u_1(t)\bigl(1-e^{-(t-1)}\bigr).\]
\end{frame}

% 29
% ANIM: system_response
\visualframe{המחשה: קלט, תגובת אימפולס ופלט}{12_system.png}{הפלט אינו זהה לקלט: הוא תוצאה של צבירת תגובת המערכת לאורך הזמן.}

% 30
\begin{frame}{התנהגות ארוכת טווח ומצב מתמיד}
כדוגמה נבחן מערכת מסדר ראשון:
\[y'(t)+y(t)=\sin t,\qquad y(0)=0.\]
פתרונה הוא
\bigformula{y(t)=\underbrace{\tfrac12 e^{-t}}_{\text{חולף}}
+\underbrace{\tfrac12(\sin t-\cos t)}_{\text{מתמיד}}}
\begin{spaceditems}
\item הרכיב החולף דועך עם הזמן
\item הרכיב המתמיד ממשיך לתנוד בתדירות הכוח החיצוני
\item תחום $s$ מסייע להבחין בין תרומות מסוגים שונים לפתרון
\end{spaceditems}
\end{frame}

% 31
\visualframe{המחשה: תגובה חולפת ותגובה מתמידה}{13_steady.png}{לאחר שהרכיב החולף נחלש, ההתנהגות נשלטת על ידי התגובה המתמידה.}

% 32
\begin{frame}{סיכום: מה למדנו על התמרת לפלס?}
\begin{spaceditems}
\item \highlight{הגדרה:} אינטגרל משוקלל שמעביר פונקציות לתחום מרוכב
\item \highlight{תנאי קיום:} רציפות למקוטעין וסדר אקספוננציאלי מספיקים בחצי מישור מתאים
\item \highlight{פישוט:} נגזרות והזזות מקבלות ייצוג אלגברי
\item \highlight{יישומים:} פתרון בעיות התחלה, קלט מושהה, אימפולסים וקונבולוציה
\item \highlight{רעיון מאחד:} חקר דינמיקה באמצעות התנהגותם של אקספוננטים
\end{spaceditems}
\vspace{6mm}
\begin{center}{\Large\color{cyan}משוואות דיפרנציאליות $\longrightarrow$ אלגברה בתחום $s$}\end{center}
\end{frame}

% 33
\begin{frame}{מקורות ותודה}
\begin{itemize}
\setlength{\itemsep}{10pt}
\item \textenglish{Joel L. Schiff}, \textenglish{The Laplace Transform: Theory and Applications}, Springer
\item \textenglish{3Blue1Brown}, סרטוני ההסבר על אקספוננטים מרוכבים והתמרת לפלס — השראה רעיונית להמחשות
\item ההמחשות והסרטונים המצורפים נבנו כתוצרים מקוריים של הפרויקט
\end{itemize}
\vspace{9mm}
\begin{center}
{\Huge\color{cyan}תודה על ההקשבה}\\[4mm]
{\LARGE\color{gold}שאלות?}
\end{center}
\end{frame}

\end{document}
